\begin{minipage}{\linewidth}
\subsubsection{Spezialfall: Regel von Sarrus}
\label{subsubsec:linear_systems_sarrus_rule}

Für $3\times3$-Determinanten existiert neben der Laplace-Entwicklung ein
spezielles Rechenschema. Dieses ist häufig schneller anzuwenden, lässt sich
jedoch nicht auf Determinanten höherer Ordnung übertragen.

\begin{KeyPointBox}{Regel von Sarrus}{linear_systems_sarrus_rule}
  Im Spezialfall einer $3\times3$-Matrix lässt sich die Determinante auch mit
  der \emph{Regel von Sarrus} \EnglishTerm{Sarrus' rule} berechnen:

  \begin{center}
    \CachedTikzFigure{content/resources/chapters/linear_systems/tikz/sarrus_rule.tex}
  \end{center}
  \vspace{-0.3\baselineskip}
  \begin{align*}
    \det(A)
    ={}&\textcolor{THMBlue}{a_{11}a_{22}a_{33}
       +a_{12}a_{23}a_{31}
       +a_{13}a_{21}a_{32}}\\
     &\textcolor{red!80!black}{-a_{13}a_{22}a_{31}
       -a_{11}a_{23}a_{32}
       -a_{12}a_{21}a_{33}}.
  \end{align*}
  \textbf{Achtung:} Die Regel von Sarrus gilt ausschließlich für
  $3\times3$-Matrizen.
\end{KeyPointBox}
\end{minipage}

\IfClozeTF{\clearpage}{}
\begin{ExampleBox}{Determinante mit der Regel von Sarrus}{linear_systems_sarrus_example}
  Wir berechnen die Determinante aus
  Beispiel~\ref{ex:linear_systems_three_by_three_determinant} erneut:
  \[
    D=\det\!\begin{pmatrix}
      1&2&3\\
      4&5&6\\
      7&8&0
    \end{pmatrix}.
  \]
  \noindent\textbf{Sarrus-Schema (Nebenrechnung):}
  \[
    \begin{array}{|ccc|cc}
      1&2&3&\textcolor{green!50!black}{1}&\textcolor{green!50!black}{2}\\
      4&5&6&\textcolor{green!50!black}{4}&\textcolor{green!50!black}{5}\\
      7&8&0&\textcolor{green!50!black}{7}&\textcolor{green!50!black}{8}
    \end{array}
  \]
  \begin{ClozeArea}[32]{7}
    \[
      \begin{aligned}
        D
        &=\textcolor{THMBlue}{1\cdot5\cdot0+2\cdot6\cdot7+3\cdot4\cdot8}\\
        &\phantom{={}}\textcolor{red!80!black}{-3\cdot5\cdot7-1\cdot6\cdot8-2\cdot4\cdot0}\\[0.6em]
        &=0+84+96-105-48-0=27.
      \end{aligned}
    \]
  \end{ClozeArea}
\end{ExampleBox}
